%% SimulacraBench Technical Report -- team Vecture
%% Built from the organizers' template (template_original.tex); section skeleton unchanged.
\documentclass{article}
%% The public neurips_2026.sty defines the footer track name only for its track
%% options; the competition template passes [final] alone, so define it here.
\makeatletter
\providecommand{\@trackname}{\@neuripsordinal\ Conference on Neural Information Processing Systems (NeurIPS \@neuripsyear). Competition Track: SimulacraBench.}
\makeatother
\usepackage[final]{neurips_2026}

\usepackage[utf8]{inputenc}
\usepackage[T1]{fontenc}
\usepackage{hyperref}
\usepackage{url}
\usepackage{booktabs}
\usepackage{amsfonts}
\usepackage{amsmath}
\usepackage{amssymb}
\usepackage{nicefrac}
\usepackage{microtype}
\usepackage{xcolor}
\usepackage{graphicx}
\usepackage{enumitem}

%% Values filled at release time.
\newcommand{\authoremail}{metin@vecture.be}
\newcommand{\coderepo}{https://github.com/mhatipog/simulacrabench-vecture}
\newcommand{\codecommit}{c5cefd6651df}


\newcommand{\teamid}{Vecture}
\newcommand{\cbox}{$\square$}
\newcommand{\cboxchecked}{$\boxtimes$}

\title{SimulacraBench Technical Report:\\ Cross-fitted, Gate-factorised Stacking of Small Survey Experts}

\author{%
  Metin Hatipogullari\\
  Vecture BV, Belgium\\
  \texttt{\authoremail}
}

\begin{document}

\maketitle

\begin{center}
\begin{tabular}{l}
\toprule
\textbf{Team ID:} \teamid \\
\bottomrule
\end{tabular}
\end{center}

\begin{abstract}
We predict a calibrated distribution over every withheld answer of a held-out respondent from that respondent's background (\texttt{GIVEN}) fields. Inside the evaluation container we fit a small, diverse set of experts (LightGBM, CatBoost, a multitask embedding network, empirical-Bayes tables, a multinomial logit and two kernel nearest-neighbour estimators), cross-fit all of them, and stack them per question on out-of-fold predictions. For questions behind skip logic the log score factorises exactly into a ``was it asked'' term and an ``answer if asked'' term; we stack the two separately, each on the rows that inform it, and recalibrate the skip probability by a smoothed isotonic map that is accepted only if it wins a two-way cross-check. A pretrained tabular in-context model (Kumo) enters as a validated blend on UNICEF and World Bank and as a capped, failure-tolerant expert on UNHCR. A budget controller times every expert on the host and chooses folds and experts to fit the wall clock. Test phase (leaderboard, two decimals): skill 0.49 (UNICEF 0.34, World Bank 0.24, UNHCR 0.88). Paired ablations on a public survey and on synthetic mechanism tests attribute the gains to stacking with per-question weights and to the factorised gate.
\end{abstract}

%% ---------------------------------------------------------------------------
\section{Method summary}
\label{sec:method}

\subsection{Method type}
Components are described in Section~\ref{sec:pipeline}.

\begin{table}[h]
  \centering
  \small
  \begin{tabular}{@{}clcl@{}}
    \toprule
    \cbox & LLM-simulated respondents (prompting, personas) & \cbox & Fine-tuned language model \\
    \cbox & Tabular / conditional generative model (e.g.\ VAE, diffusion, GAN, CTGAN) & \cboxchecked & Supervised prediction (GBDT, neural net, linear/GLM) \\
    \cboxchecked & Bayesian / hierarchical / multilevel model & \cbox & Small-area estimation or post-stratification (e.g.\ MRP) \\
    \cbox & Reweighting, raking, or calibration of the released sample & \cboxchecked & Imputation / matching / nearest-neighbour \\
    \cbox & Retrieval over released or external microdata & \cbox & Agent-based or mechanistic simulation \\
    \cboxchecked & Ensemble or stacking of multiple components & \cboxchecked & Other: pretrained tabular in-context model \\
    \bottomrule
  \end{tabular}
\end{table}

\subsection{Pipeline}
\label{sec:pipeline}
One \texttt{predict(frame, schema)} call per instrument. (i) Encode \texttt{GIVEN} fields; take labels from visible rows only. (ii) Probe every expert's cost on this host and pick the largest plan (expert set, $K\in\{2,\dots,5\}$ folds) that fits this call's share of the wall clock. (iii) Cross-fit: fit each expert $K$ times and predict its held-out fold, giving out-of-fold (OOF) predictions for every visible row, then once on all visible rows to predict the scored rows. (iv) Stack per question on the OOF predictions; for gated questions stack the gate and the answer separately (\S\ref{sec:agg}). (v) Recalibrate gate probabilities (\S\ref{sec:post}). (vi) Kumo: on UNICEF and World Bank a validated mixture with the stack; on UNHCR a cross-fitted expert inside the stack, time-capped, whose failure is logged and leaves the stack unchanged. Components marked above: \emph{supervised} = LightGBM~\cite{ke2017lightgbm}, CatBoost~\cite{prokhorenkova2018catboost}, network, logit; \emph{Bayesian} = empirical-Bayes smoothed tables; \emph{nearest-neighbour} = kernel matching on the background block; \emph{stacking} = constrained per-question stacking~\cite{wolpert1992stacked,breiman1996stacked} with cross-fitting~\cite{chernozhukov2018double}; \emph{other} = Kumo~\cite{kumo2026}.

\subsection{Inputs and feature construction}
Only \texttt{GIVEN} fields are features: every scored respondent has them and nothing else, so answers to one question cannot inform another at prediction time. Fields are integer-coded from the schema's option list; banded fields (``6-10'', ``25-34'') are also given as numbers to LightGBM, plus totals over count fields that share an option list (household roster). Question text, weights and row identifiers are not used.

\subsection{Model or simulator}
Experts, each producing a distribution over a question's options (including the final ``not asked'' sentinel for gated questions): \emph{direct} LightGBM and CatBoost multiclass models per question; a multitask network (field embeddings, two 256-unit layers, one softmax head per question plus a gate head); a multinomial logit for all questions in one linear map; empirical-Bayes tables choosing between a one-field table, a coarse two-field table and a tempered naive-Bayes combination on an inner split; two kernel neighbour estimators weighting a visible row by $\beta^{d}$, $d$ = number of differing \texttt{GIVEN} fields, $\beta\in\{0.1,0.3\}$, smoothed to the marginal. \emph{Gate} experts: one binary LightGBM per distinct skip pattern. \emph{Answer-if-asked} experts: LightGBM and logit fitted on asked rows only. Questions are treated as conditionally independent given \texttt{GIVEN}, which is exact for this task's per-cell score. Hyperparameters are in Appendix~\ref{app:config}.

\subsection{Fitting, training, or prompting}
Everything is fitted inside the container on the visible rows of the frame it receives (TRAIN in Development, TRAIN+DEV in Test); no model is trained offline except Kumo, used as released. Data-chosen settings (table type, network epochs by early stopping on a 12\% split) are chosen inside training folds. Development of the method used the Development leaderboard and the public proxy data of Table~\ref{tab:external}; no hidden labels or leaderboard probing were used.

\subsection{Prediction and aggregation}
\label{sec:agg}
For an ungated question, $p=\sum_m w_m p_m + w_0\,\pi$ with $\pi$ the visible marginal and $w$ on the simplex. For a gated question with sentinel probability $g$ and answer distribution $c$, the cell log loss is $-\log g$ if skipped and $-\log(1-g)-\log c_y$ if asked; the two terms share no parameters, so the gate stack (every expert's sentinel probability, the binary gate model, the marginal skip rate) is fitted on all rows and the answer stack (renormalised non-sentinel parts, asked-only experts, asked marginal) on asked rows only. Weights minimise the floored OOF log loss by SLSQP with a quadratic pull toward weights pooled over the instrument (penalty $0.004$, $0.02$ for answer stacks with $\le 400$ asked rows); a temperature $T\in\{0.8,\dots,1.5\}$ is then chosen per question.

\subsection{Post-processing and calibration (if applicable)}
\label{sec:post}
A convex blend cannot be sharper than its sharpest member, while structurally determined skips deserve near-certain gates. We fit an isotonic map from the stacked OOF gate probability to the observed skip indicator; each block of $n$ rows with mean input $x$ and rate $r$ is smoothed to $(nr+0.5+x)/(n+2)$, so it is never 0 or 1. The map is used for a question only if it lowers the loss in a two-way cross-check on disjoint halves of the OOF rows. Declared gates are never imposed: the organizers report gate/data agreement as low as 0.29 for UNHCR, and an earlier entry of ours that forced declared gates fell from 0.85 to 0.56 UNHCR skill in Development. Outputs are clipped at $10^{-9}$ and renormalised.

%% ---------------------------------------------------------------------------
\section{Models, data, and compute}
\label{sec:resources}

\paragraph{Pretrained models and off-the-shelf components.} Table~\ref{tab:models} lists every component with non-trivial modelling content.
\begin{table}[h]
  \centering
  \small
  \caption{Pretrained models and off-the-shelf components.}
  \label{tab:models}
  \begin{tabular}{lll}
    \toprule
    Component & Version / checkpoint & Role in pipeline  \\
    \midrule
    \texttt{nvidia/Kumo-Tabular} via \texttt{structured-data-models} & 0.2.1; \texttt{small}, \texttt{large} & in-context expert (\S\ref{sec:pipeline}) \\
    LightGBM & 4.6.0 & direct, gate, answer experts \\
    CatBoost & 1.2.8 & direct expert \\
    PyTorch, SciPy (SLSQP), NumPy & container versions & network, logit, stacking \\
    \bottomrule
  \end{tabular}
\end{table}

\paragraph{Data.} The SimulacraBench TRAIN rows visible in each phase were the only data the submission saw. Development used the Development leaderboard and organizer sandboxes; the public survey in Table~\ref{tab:external} was used only offline to validate design choices.

\begin{table}[h]
  \centering
  \small
  \caption{External datasets used.}
  \label{tab:external}
  \begin{tabular}{p{2.6cm}p{4.8cm}p{1.8cm}p{3.4cm}}
    \toprule
    Dataset & Source / URL & Size & Use in pipeline \\
    \midrule
    General Social Survey 2022 (NORC) & \url{https://gss.norc.org} (public Stata release) & 4,149 resp. & offline validation only; not in the submission \\
    \bottomrule
  \end{tabular}
\end{table}

\paragraph{Compute.} Hosted: one H100, 8 CPUs, 16\,GB, no network; Development limit 900\,s, Test 3,600\,s for all three instruments. Our Development run of this exact archive took 49.1\,s. All local work used one 16-thread CPU workstation; no GPU.

%% ---------------------------------------------------------------------------
\section{Results}
\label{sec:results}

\subsection{Test-phase leaderboard results}
\begin{table}[h]
  \centering
  \small
  \caption{Official SimulacraBench results for team \teamid{} (submission archive \texttt{SimulacraBench\_F5\_KumoStackRobust.zip}, SHA-256 \texttt{07e4ba37\ldots c76f75}). Test values are as displayed on the leaderboard (two decimals). Development values are from the scoring output and are noised and bucketed by design.}
  \label{tab:results}
  \begin{tabular}{lrrrr}
    \toprule
    Split & Skill (Grand Prize) & UNICEF & World Bank & UNHCR \\
    \midrule
    Development (13 Nov 2026) & 0.4878 & 0.3499 & 0.2486 & 0.8650 \\
    Test (submission 1053704, 15 Nov 2026) & 0.49 & 0.34 & 0.24 & 0.88 \\
    Organizer baseline (crowd marginal; 0 by definition) & 0.00 & 0.00 & 0.00 & 0.00 \\
    \bottomrule
  \end{tabular}
\end{table}

\subsection{Breakdown by population, survey, or question type}
The leaderboard reports one skill per survey. Our weakest surveys are World Bank (0.24) and UNICEF (0.34), which are also flat across teams: in the snapshot of 17 Nov 2026 (46 entries), 38 entries show UNICEF 0.34 and none higher, and 18 show World Bank 0.24, none higher. Differences between leading entries are almost entirely on UNHCR, where 16 entries show 0.88 or 0.89. In Development, UNHCR moved up one bucket when the kernel neighbour experts were added, and World Bank flipped between adjacent buckets with byte-identical code (Appendix~B).

\subsection{Ablations on proxy data (not leaderboard numbers)}
\label{sec:ablation}
\begin{table}[t]
\centering
\scriptsize
\setlength{\tabcolsep}{3pt}
\caption{Paired ablations (Kumo off, identical expert fits): change in mean log score per cell, milli-nats, when the component is removed (negative = it helps; bold: 95\% bootstrap interval excludes 0). GSS: real survey; Gate/Twin: synthetic (Appendix~B).}
\label{tab:ablation}
\begin{tabular}{lrrrrrr}
\toprule
Variant & GSS-1 & GSS-2 & GSS-3 & Gate-A & Twin & Twin-HH \\
\midrule
Full stack (submitted), skill & 0.4209 & 0.4745 & 0.4436 & 0.2242 & 0.2872 & 0.1316 \\
\midrule
$-$ LightGBM (direct) & -0.0 & -0.2 & \textbf{-0.3} & -0.0 & \textbf{-0.1} & -0.0 \\
$-$ CatBoost & \textbf{-0.9} & \textbf{-1.5} & \textbf{-0.5} & -0.4 & -0.3 & -0.2 \\
$-$ multitask net & -0.8 & -0.5 & +0.1 & \textbf{-2.8} & \textbf{-6.2} & +0.5 \\
$-$ empirical-Bayes tables & +0.2 & +0.1 & +0.2 & \textbf{-0.5} & \textbf{-0.3} & -0.1 \\
$-$ multinomial logit (direct+cond.) & -0.0 & +0.1 & -0.1 & \textbf{-2.7} & \textbf{-1.2} & \textbf{-1.6} \\
$-$ kernel neighbours & \textbf{+0.2} & +0.1 & +0.2 & -0.0 & \textbf{-82.3} & +1.7 \\
$-$ gate/answer-if-asked LightGBM & \textbf{-0.5} & \textbf{-0.4} & \textbf{-0.7} & +0.0 & +0.0 & \textbf{+0.4} \\
no gate factorisation (one stack per item) & \textbf{-4.3} & \textbf{-4.5} & \textbf{-3.7} & \textbf{-0.8} & \textbf{-9.1} & \textbf{-4.6} \\
no isotonic gate recalibration & \textbf{-0.4} & \textbf{-0.5} & \textbf{-0.4} & \textbf{+0.3} & \textbf{-0.8} & \textbf{-0.7} \\
no temperature & \textbf{-0.3} & \textbf{-0.2} & +0.0 & -0.0 & \textbf{-3.4} & \textbf{-1.0} \\
pooled weights only (no per-item) & \textbf{-1.4} & \textbf{-1.0} & -0.1 & -0.5 & \textbf{-12.6} & -1.2 \\
stack fitted on 950 rows (V81-style) & \textbf{-1.1} & \textbf{-2.3} & \textbf{-1.0} & \textbf{-1.0} & \textbf{-3.5} & \textbf{-1.3} \\
equal-weight average of experts & \textbf{-58.6} & \textbf{-65.2} & \textbf{-64.9} & \textbf{-38.6} & \textbf{-51.5} & +0.4 \\
single best expert (selected out-of-fold) & \textbf{-13.4} & \textbf{-12.6} & \textbf{-11.4} & \textbf{-18.7} & \textbf{-101.2} & -0.6 \\
\bottomrule
\end{tabular}
\end{table}

The full variant reproduces the submitted predictions exactly; refitting all experts with two other seeds moves the full stack by at most 0.11 and flips no non-null effect (Appendix~B). On the real survey, removing the gate factorisation is the largest single loss (3.7--4.5 milli-nats per cell), and it hurts on every set. Fitting the stack on 950 rows, as our earlier ensemble did, rather than on all out-of-fold rows costs 1.0--2.3 there and 0.9--3.5 elsewhere; per-question weights and isotonic gate recalibration add 0.1--1.4 and 0.4--0.5 on real skip patterns. The stack beats its best single expert (selected out-of-fold) by 11--13 and an equal-weight average by about 60. Among the direct experts only CatBoost is indispensable on all three GSS studies (0.6--1.5); the dedicated gate and answer-if-asked models add 0.4--0.7. The kernel neighbours are worth 82 milli-nats when scored rows have visible cross-round twins (Twin) and nothing on GSS; when scored rows never have twins but visible rows do (Twin-HH), removing them helps slightly (+1.7, interval includes 0), consistent with out-of-fold evidence overstating them there. Against our earlier four-expert ensemble the submitted stack gains 6.3--10.6 milli-nats per cell on GSS and 3.1--5.6 on the gate simulators, all intervals excluding zero.


%% ---------------------------------------------------------------------------
\section{Reproducibility}
\label{sec:repro}

\begin{itemize}[leftmargin=*]
  \item Code repository: \url{\coderepo} (licence: MIT).
  \item Commit used for the test-phase submission: \texttt{\codecommit}; the archive's files are in \texttt{submission/}.
  \item Reproduce the submission file: \texttt{python tools/build\_submission\_zip.py} (byte-identical archive, SHA-256 checked). The platform runs its \texttt{predict()} on the released data; locally: \texttt{python eval/run\_eval.py submission <data> 2}.
  \item Environment: Python 3.13; \texttt{submission/requirements.txt} pins LightGBM 4.6.0, CatBoost 1.2.8, \texttt{structured-data-models} 0.2.1; PyTorch, SciPy from the hosted image. Weights: \texttt{nvidia/Kumo-Tabular} (public on Hugging Face).
  \item Ablations and tables: \texttt{eval/ablate\_all.sh}, \texttt{eval/make\_tables.py}.
\end{itemize}

%% ---------------------------------------------------------------------------
\section{Limitations}
\label{sec:limitations}

The method sees only background fields, so skipping driven by unobserved answers stays uncertain. Whether Kumo covered any UNHCR question in the hosted test run is not observable (no logs are returned); our ablations run with Kumo off on CPU. The time-adaptive plan and GPU nondeterminism make reruns hardware-dependent. Proxy evidence comes from one public survey family and synthetic generators.

%% ---------------------------------------------------------------------------
\bibliographystyle{plainnat}
\bibliography{refs}

%% ---------------------------------------------------------------------------
\clearpage
\appendix

\section{Full configuration}
\label{app:config}
No prompts are used. All settings below are fixed in \texttt{submission/} (file in brackets); nothing is tuned per instrument except where a selection rule is stated. Seed 20261111 throughout.

\paragraph{Features (\texttt{sb\_core.py}).} \texttt{GIVEN} fields integer-coded (0 = unseen level). Numeric companions for fields whose levels are mostly numbers or bands: band midpoint (``25-34''$\to$29.5), lower bound for open bands (``11+''$\to$11); plus the sum over three or more such fields sharing one level list (household roster) and the difference for exactly two. LightGBM: codes as categoricals plus numeric companions; CatBoost: codes as categoricals; logit and kernel experts: one-hot codes; network: 14-dimensional embeddings per field.

\paragraph{Experts (\texttt{sb\_experts.py}, \texttt{sb\_eb.py}, \texttt{survey\_experts.py}).}
\begin{itemize}[leftmargin=*,itemsep=1pt]
\item LightGBM direct and answer-if-asked: 96 trees, learning rate 0.055, 11 leaves, \texttt{min\_child\_samples} $\min(65,n/12)$ (at least 5), $\lambda_2=15$, $\lambda_1=0.18$, column sample 0.95, 63 bins; output mixed $0.99\,p+0.01\,\pi$. Gate model (one per distinct skip pattern): 160 trees, rate 0.07, 15 leaves, \texttt{min\_child\_samples} 20, $\lambda_2=2$, mixed with $0.002\,\pi$.
\item CatBoost: depth 5, rate 0.058, $\ell_2=7$, random strength 0.8, Bernoulli subsample 0.85, one-hot up to 12 levels; 115 iterations (86 when an instrument has more than 120 questions); mixed with $0.01\,\pi$.
\item Network: embeddings (dim 14) $\to$ 256 (GELU, LayerNorm, dropout 0.15) $\to$ 256 (GELU, dropout 0.12) $\to$ per-question softmax heads and a gate head (gated questions factorise sentinel and answer); AdamW, rate 0.002, weight decay 0.065, batch 384, gradient clip 2; epochs chosen once by early stopping (patience 5, at most 20) on a 12\% split of the first training fold, then reused.
\item Multinomial logit: one linear map for all direct and answer-if-asked heads, ridge 0.8$/n$, Adam rate 0.07 with cosine decay, 140 steps, biases initialised at log marginals.
\item Empirical-Bayes tables: Dirichlet-smoothed $P(y\mid x_i)$ (mass 22) or $P(y\mid x_i,x_j)$ (mass 42, pairs with at most 360 cells), or naive-Bayes products of the best 2, 3 or 5 single tables at temperature 0.25/0.5/0.75, or the marginal; chosen per question on an inner 80/20 split by mean advantage minus $0.65\times$ standard error minus a complexity cost.
\item Kernel neighbours: weight $\beta^{d}$ with $\beta\in\{0.1,0.3\}$, $d$ = differing \texttt{GIVEN} fields, plus 0.02 pseudo-rows of the marginal.
\item Kumo (\texttt{kumo\_parents.py}): UNICEF all 12 questions, World Bank questions with at most 12 options; Kumo Small, 2 estimators, context = visible rows outside a validation subset of $\min(950,\max(45,0.165n))$ rows. Mixture weight $\alpha\in\{0,\dots,1\}$ and temperature $\in\{0.85,\dots,1.65\}$ chosen per question family on that subset with penalties $0.0012\alpha^2$ and $0.0005|\log T|$, then per question shrunk toward the family value. UNHCR: questions with at most 12 options, ordered by marginal entropy; Kumo Large if its first fold projects $\ge 90\%$ coverage within the cap, else Small; fallback ladder (Small, float32, no numeric extras); 4 folds plus a final pass; capped at 45\% of the UNHCR time share; enters the stack as an expert (marginal where not covered).
\end{itemize}

\paragraph{Plan and folds (\texttt{main.py}).} Budget = $0.8\times$ phase limit minus elapsed time, shared across the three calls by fixed work weights (12 questions: 1.5, 76: 3.5, 213: 5.0). Candidate expert sets in order: all six; without CatBoost; without CatBoost and logit; without CatBoost and network; LightGBM, logit, tables; logit and tables. For each, the largest $K\in\{5,4,3,2\}$ with $(K+1)\times$ probed cost within budget. Folds are random (seed $+J$).

\paragraph{Stacking and calibration (\texttt{sb\_stack.py}).} SLSQP on the simplex, at most 80 iterations; pooled weights from up to 60,000 stacked OOF rows with pull 0.003 toward uniform; per-question pull 0.004 toward pooled (0.02 for answer stacks with $\le 400$ asked rows). Temperature grid $\{0.8,0.9,1,1.1,1.25,1.5\}$ with penalty $4\cdot 10^{-4}\log^2 T$. Isotonic gate map on at most 300 quantile bins, pool-adjacent-violators, block smoothing $(nr+0.5+x)/(n+2)$, accepted if two-way cross-check gain $>2\cdot10^{-4}$ nats per row; requires $\ge 200$ rows with both outcomes present. Instruments with fewer than 90 visible rows use a similarity-weighted Dirichlet estimate instead (never triggered on the official data).


\section{(Optional) Additional results}
These are offline proxy results, not leaderboard numbers. Datasets: GSS-1..3, three studies of 70 questions each from the 2022 General Social Survey (3,000 visible, 600 scored respondents, real skip patterns); Gate-A/B/C, synthetic data on the UNHCR schema whose gates are obeyed at rates 0.29--0.98, with 50\%, 0\% and 50\% of gates determined by background fields (C: the remaining gates obeyed exactly); Twin and Twin-HH, synthetic data in which 39\% of rows re-interview an earlier household in a second round, with roles assigned per row (Twin) or per household (Twin-HH); SB-*, organizer sandboxes at the declared sizes. Builders and seeds are in \texttt{eval/make\_datasets.sh}.

\begin{table}[t]
\centering
\tiny
\setlength{\tabcolsep}{3pt}
\caption{All ablations on all eleven evaluation sets (milli-nats per cell, variant minus full stack; bold: 95\% interval excludes zero). Intervals are listed in \texttt{results/ablation\_full.csv}.}
\label{tab:ablation-all}
\begin{tabular}{lrrrrrrrrrrr}
\toprule
Variant & GSS-1 & GSS-2 & GSS-3 & Gate-A & Gate-B & Gate-C & Twin & Twin-HH & SB-UNICEF & SB-WB & SB-UNHCR \\
\midrule
Full stack (submitted), skill & 0.4209 & 0.4745 & 0.4436 & 0.2242 & 0.1208 & 0.2455 & 0.2872 & 0.1316 & 0.4025 & 0.2809 & 0.2308 \\
\midrule
$-$ LightGBM (direct) & -0.0 & -0.2 & \textbf{-0.3} & -0.0 & -0.2 & +0.1 & \textbf{-0.1} & -0.0 & +0.0 & +0.0 & -0.0 \\
$-$ CatBoost & \textbf{-0.9} & \textbf{-1.5} & \textbf{-0.5} & -0.4 & \textbf{-1.0} & \textbf{-0.7} & -0.3 & -0.2 & +0.0 & -0.4 & \textbf{-0.5} \\
$-$ multitask net & -0.8 & -0.5 & +0.1 & \textbf{-2.8} & \textbf{-3.0} & -1.0 & \textbf{-6.2} & +0.5 & -0.1 & \textbf{-3.9} & -0.6 \\
$-$ empirical-Bayes tables & +0.2 & +0.1 & +0.2 & \textbf{-0.5} & \textbf{-1.0} & \textbf{-1.0} & \textbf{-0.3} & -0.1 & -0.2 & \textbf{-0.3} & +0.0 \\
$-$ multinomial logit (direct+cond.) & -0.0 & +0.1 & -0.1 & \textbf{-2.7} & \textbf{-3.1} & \textbf{-1.5} & \textbf{-1.2} & \textbf{-1.6} & \textbf{-0.5} & +0.0 & +0.0 \\
$-$ kernel neighbours & \textbf{+0.2} & +0.1 & +0.2 & -0.0 & +0.0 & +0.0 & \textbf{-82.3} & +1.7 & -0.0 & -0.0 & -0.1 \\
$-$ gate/answer-if-asked LightGBM & \textbf{-0.5} & \textbf{-0.4} & \textbf{-0.7} & +0.0 & -0.2 & -0.3 & +0.0 & \textbf{+0.4} & +0.1 & \textbf{-0.3} & -0.0 \\
no gate factorisation (one stack per item) & \textbf{-4.3} & \textbf{-4.5} & \textbf{-3.7} & \textbf{-0.8} & \textbf{-1.0} & \textbf{-1.0} & \textbf{-9.1} & \textbf{-4.6} & \textbf{-1.5} & \textbf{-2.0} & \textbf{-0.4} \\
no isotonic gate recalibration & \textbf{-0.4} & \textbf{-0.5} & \textbf{-0.4} & \textbf{+0.3} & -0.1 & +0.0 & \textbf{-0.8} & \textbf{-0.7} & \textbf{-0.4} & \textbf{-0.3} & +0.0 \\
no temperature & \textbf{-0.3} & \textbf{-0.2} & +0.0 & -0.0 & -0.2 & -0.2 & \textbf{-3.4} & \textbf{-1.0} & +0.0 & -0.1 & -0.0 \\
pooled weights only (no per-item) & \textbf{-1.4} & \textbf{-1.0} & -0.1 & -0.5 & -0.2 & -0.4 & \textbf{-12.6} & -1.2 & \textbf{-0.4} & -0.4 & +0.2 \\
stack fitted on 950 rows (V81-style) & \textbf{-1.1} & \textbf{-2.3} & \textbf{-1.0} & \textbf{-1.0} & \textbf{-0.9} & \textbf{-1.8} & \textbf{-3.5} & \textbf{-1.3} & \textbf{-2.2} & \textbf{-1.2} & -1.3 \\
equal-weight average of experts & \textbf{-58.6} & \textbf{-65.2} & \textbf{-64.9} & \textbf{-38.6} & \textbf{-19.2} & \textbf{-38.2} & \textbf{-51.5} & +0.4 & \textbf{-21.1} & \textbf{-18.3} & \textbf{-6.5} \\
single best expert (selected out-of-fold) & \textbf{-13.4} & \textbf{-12.6} & \textbf{-11.4} & \textbf{-18.7} & \textbf{-22.7} & \textbf{-13.2} & \textbf{-101.2} & -0.6 & \textbf{-5.5} & \textbf{-6.8} & \textbf{-2.8} \\
\bottomrule
\end{tabular}
\end{table}


\begin{table}[t]
\centering\small
\caption{Earlier ensemble (V81) against the submitted stack (F5, Kumo off, full plan), paired on identical cells. Mean log loss per cell (nats) and gain of F5 with 95\% respondent-bootstrap intervals.}
\label{tab:v81f5}
\begin{tabular}{lrrl}
\toprule
Dataset & V81 loss & F5 loss & Gain (95\% CI) \\
\midrule
GSS-1 & 0.5842 & 0.5754 & $+0.0089$ [$+0.0072$, $+0.0104$] \\
GSS-2 & 0.5184 & 0.5078 & $+0.0106$ [$+0.0086$, $+0.0126$] \\
GSS-3 & 0.5616 & 0.5553 & $+0.0063$ [$+0.0044$, $+0.0081$] \\
Gate-A & 0.7456 & 0.7407 & $+0.0049$ [$+0.0029$, $+0.0071$] \\
Gate-B & 0.8339 & 0.8308 & $+0.0031$ [$+0.0007$, $+0.0055$] \\
Gate-C & 0.6713 & 0.6657 & $+0.0056$ [$+0.0038$, $+0.0077$] \\
\bottomrule
\end{tabular}
\end{table}


\begin{table}[h]
\centering\small
\caption{Development-phase results of this line (scoring files returned by the platform; noised and bucketed). One displayed step is 0.0156, 0.0202 and 0.0126 nats per cell on UNICEF, World Bank and UNHCR. F4 and F5 scored exactly as F3.}
\label{tab:devhist}
\begin{tabular}{llrrrrr}
\toprule
Entry & Change & Overall & UNHCR & UNICEF & World Bank & Time \\
\midrule
F1 (2 archives) & cross-fitted factorised stack & 0.4764 / 0.4814 & 0.8457 & 0.3499 & 0.2335 / 0.2486 & 30 / 53 s \\
F3 / F4 / F5 & + kernel neighbours, Kumo on UNHCR & 0.4878 & 0.8650 & 0.3499 & 0.2486 & 34--49 s \\
F6 & + twin experts & 0.4764 & 0.8457 & 0.3499 & 0.2335 & 23 s \\
G1 & + kernel family & 0.4814 & 0.8457 & 0.3499 & 0.2486 & 45 s \\
H1 & + household-grouped folds & 0.4828 & 0.8650 & 0.3499 & 0.2335 & 25 s \\
\bottomrule
\end{tabular}
\end{table}

The test submission is the F5 archive unchanged: it shares the best Development score with F3 and F4 (identical to every digit), and, unlike F6 and G1, it adds no cross-round matching; H1's paired gain was shown only on a simulator built for that mechanism.

\begin{table}[h]
\centering\small
\caption{Expert-fit seed check on GSS-1: all experts and folds refitted with three seeds, then re-stacked. Change in mean log score per cell (milli-nats) when the component is removed. Full-stack loss: 0.57535, 0.57536, 0.57547 nats.}
\label{tab:seeds}
\begin{tabular}{lrrr}
\toprule
Variant & Seed 20261111 (submitted) & Seed 1 & Seed 2 \\
\midrule
no gate factorisation & $-4.31$ & $-4.70$ & $-4.63$ \\
stack fitted on 950 rows & $-1.13$ & $-2.04$ & $-1.86$ \\
pooled weights only & $-1.37$ & $-1.86$ & $-1.61$ \\
$-$ CatBoost & $-0.94$ & $-1.02$ & $-0.98$ \\
$-$ multitask net & $-0.80$ & $-0.96$ & $-0.67$ \\
$-$ gate/answer-if-asked LightGBM & $-0.48$ & $-0.33$ & $-0.50$ \\
no isotonic gate recalibration & $-0.43$ & $-0.63$ & $-0.41$ \\
no temperature & $-0.27$ & $-0.31$ & $-0.33$ \\
$-$ kernel neighbours & $+0.17$ & $+0.16$ & $+0.15$ \\
$-$ empirical-Bayes tables & $+0.23$ & $-0.09$ & $+0.36$ \\
single best expert & $-13.36$ & $-13.21$ & $-13.30$ \\
\bottomrule
\end{tabular}
\end{table}



\end{document}
